% Requires amsmath, amssymb, amsthm and theorem/lemma/remark environments.
\section{Polynomial mistakes for fresh uniform ties and minimum-class majority}

Let $\mathcal C_j$ be a finite set of circuit descriptions of size $j$,
and write $N_j=|\mathcal C_j|$. Distinct descriptions computing the same
function are counted separately. Let $f$ have a description of size $s$.
All labels are noiseless, every label is revealed after prediction, and
all subsequent predictions must be consistent with every revealed label.
No algorithmic efficiency is asserted.

\subsection{Fresh uniform tie-breaking: arbitrary adaptive input order}

\begin{theorem}[Fresh uniform minimum-description prediction]\label{thm:uniform-min}
On each round, choose a fresh uniformly random description from the
currently consistent descriptions of minimum size, and predict with it.
The input may be chosen adaptively using the entire past and the learner's
previous predictions, but must be chosen before the fresh random choice
of the current description. Put
\[
 B_s=s+\sum_{j=0}^s\log\max\{1,N_j\}.
\]
Then the total number $M_\infty$ of mistakes over an arbitrarily long,
possibly infinite, input sequence satisfies
\[
 \mathbb E M_\infty\le B_s,
 \qquad
 \mathbb P\{M_\infty\ge(e-1)B_s+\log(1/\delta)\}\le\delta.
\]
For the all-binary-gate convention of this paper, and hence also for NAND,
\[
 B_s\le s+s(s+1)\log(4(d+s+1)),
\]
as in \eqref{eq:tier-count}. The empty size-zero class is omitted.
In particular this is a polynomial total-mistake guarantee for exact
minimum-size prediction under fully adaptive nonanticipating input order.
The learner need not know $s$.
\end{theorem}

\begin{proof}
Immediately before a label is revealed, let $j$ be the minimum surviving
size, let $n$ be the number of surviving descriptions at that size, and
let $q$ be the fraction of these descriptions wrong on the current input.
The conditional probability of the learner's mistake is exactly $q$.
If $q<1$, the surviving count at this same size becomes $n(1-q)$ and
\[
 q\le-\log(1-q)=\log n-\log(n(1-q)).
\]
Thus the sum of all such $q$ while size $j$ is the minimum is at most
$\log N_j$. If $q=1$, the size-$j$ class is exhausted and this transition
costs one. Minimum size never decreases. It never exceeds $s$, and the
size-$s$ class cannot be exhausted because the target survives. There are
therefore at most $s$ exhausted tiers. On every realization,
\[
 \sum_{t\ge1}q_t\le B_s.
\]
For the actual mistake indicators $Z_t$, conditional on the past and the
current input, $\mathbb E Z_t=q_t$ and
\[
 \mathbb E\exp\{Z_t-(e-1)q_t\}
 =e^{-(e-1)q_t}(1+(e-1)q_t)\le1.
\]
Consequently $\exp\{M_T-(e-1)\sum_{t\le T}q_t\}$ is a nonnegative
supermartingale, also when the next input is chosen from the full past.
Monotone convergence proves the expectation bound, and Ville's inequality
proves the time-uniform tail bound. The counting estimate is \eqref{eq:tier-count}.
\end{proof}

\begin{remark}[What the tie rule does and does not say]
Uniformity is over the specified finite set of minimum-size descriptions,
not automatically over distinct Boolean functions. Uniform sampling over
the latter also works if the corresponding finite counts are used.
The random choice must be fresh and private until the current input is
fixed. A single random ordering of descriptions, reused forever, is not
conditionally uniform after an adaptive input chooser has observed
predictions, so the above proof does not apply to that rule.
\end{remark}

\begin{corollary}[Minimum-tier majority]\label{cor:minimum-majority}
Predicting the majority label among the currently consistent descriptions
of minimum size, with arbitrary ties, makes at most $2B_s$ mistakes on
every input sequence. This can also be implemented as a query-dependent
choice of a minimum-size consistent circuit agreeing with that majority.
\end{corollary}

\begin{proof}
The same pathwise bound on $\sum q_t$ holds without any random prediction
rule. Each majority mistake has $q_t\ge1/2$.
\end{proof}


\subsection{The conjunction ladder under fresh uniform ties}
\begin{proposition}[Fresh uniform ties on the same ladder]
\label{prop:uniform-ties-ladder}
Use the all-binary-gate or AND/OR/NOT basis, with all gates charged.
Suppose that before every prediction the learner freshly samples
uniformly from all consistent circuits of minimum gate count.
Consider the level-by-level input order from
Section~\ref{sec:exact-occam-lower}, with the order
within each level chosen before the current fresh random draw.
For each $k=2,\ldots,d-1$, its expected mistakes in level $k$ are
exactly
\[
  H_{\binom dk}=\sum_{r=1}^{\binom dk}\frac1r.
\]
Thus its expected mistakes in these levels are
\[
  \sum_{k=2}^{d-1}H_{\binom dk}=\Theta(d^2),
\]
and the initial $d+2$ queries contribute at most $d+2$ further
mistakes.
\end{proposition}
\begin{proof}
Lemma~\ref{lem:lower-support} shows that the minimum circuits
in level $k$ compute precisely the surviving $k$-conjunctions.
Permuting input coordinates is a bijection between the size-$(k-1)$
circuits computing any two such conjunctions. They therefore have
equal numbers of minimum-size descriptions, in each of these two
bases. Since all descriptions computing a given function agree
on the entire observed history, uniformity over minimum circuit
descriptions induces uniformity over the surviving conjunctions.

If $r$ remain, any unqueried input of weight $k$ is predicted
incorrectly by exactly one of them. The conditional mistake
probability is $1/r$, and the revealed negative label removes
exactly that conjunction. Summing gives the harmonic number.
For the asymptotic estimate, $H_n\le1+\ln n$ gives an upper bound
$O(d^2)$. For $\Theta(d)$ central values of $k$,
$\binom dk=\exp(\Omega(d))$, so
$H_{\binom dk}\ge\ln\binom dk=\Omega(d)$, proving the lower bound.
\end{proof}
